
\documentclass[11pt]{article}

\usepackage[utf8]{inputenc}
\usepackage[T1]{fontenc}
\usepackage[english]{babel}
\usepackage{lmodern}
\usepackage[a4paper,margin=2.5cm]{geometry}
\usepackage{microtype}
\usepackage{booktabs}
\usepackage{tabularx}
\usepackage{array}
\usepackage{enumitem}
\usepackage{xcolor}
\usepackage{hyperref}
\usepackage[most]{tcolorbox}
\usepackage{fancyhdr}
\usepackage{titlesec}
\usepackage{caption}

\definecolor{presblue}{HTML}{17365D}
\definecolor{lightblue}{HTML}{EAF1F8}
\definecolor{softgray}{HTML}{F4F4F4}
\definecolor{decisiongreen}{HTML}{EAF4EA}
\definecolor{warningyellow}{HTML}{FFF7DF}

\hypersetup{
  colorlinks=true,
  linkcolor=presblue,
  urlcolor=presblue,
  citecolor=presblue,
  pdftitle={Spatial Intelligence for Decision Support},
  pdfauthor={SummerUW}
}

\pagestyle{fancy}
\fancyhf{}
\lhead{SummerUW}
\rhead{Spatial Intelligence}
\cfoot{\thepage}
\setlength{\headheight}{14pt}

\titleformat{\section}{\Large\bfseries\color{presblue}}{\thesection}{0.7em}{}
\titleformat{\subsection}{\large\bfseries\color{presblue}}{\thesubsection}{0.7em}{}

\newtcolorbox{executivebox}{
  colback=lightblue,
  colframe=presblue,
  title=The client situation,
  fonttitle=\bfseries,
  breakable
}

\newtcolorbox{notebookbox}{
  colback=softgray,
  colframe=black!45,
  title=Computational notebook,
  fonttitle=\bfseries,
  breakable
}

\newtcolorbox{decisionbox}{
  colback=decisiongreen,
  colframe=green!45!black,
  title=Decision supported,
  fonttitle=\bfseries,
  breakable
}

\newtcolorbox{validationbox}{
  colback=warningyellow,
  colframe=orange!55!black,
  title=For the expert to consider,
  fonttitle=\bfseries,
  breakable
}

\newcommand{\colablink}{\href{https://colab.research.google.com/}{Open the workshop notebook in Google Colab}}

\newcommand{\outputplaceholder}[1]{%
\begin{tcolorbox}[
  colback=white,
  colframe=black!30,
  title=Colab output: #1,
  fonttitle=\bfseries
]
\centering
\vspace{1.2cm}
\textit{Insert the table or figure generated by the notebook here.}
\vspace{1.2cm}
\end{tcolorbox}
}

\title{\textbf{Spatial Intelligence for Decision Support}\\[0.4em]
\large Two consulting requests for a computational tools summer school}
\author{SummerUW}
\date{August 2026}

\begin{document}

\maketitle

\begin{abstract}
This guide introduces spatial analytics through two consulting requests. In the first, a municipality needs to identify gas stations affected by a new school-distance rule and screen possible relocation areas. In the second, a public-health organization needs to distinguish different local spatial patterns of COVID-19 fatality and determine where different kinds of attention may be warranted. The guide is organized around client questions rather than GIS commands. Python produces the spatial evidence; the consultant interprets that evidence; and the client remains responsible for the final decision.
\end{abstract}

\tableofcontents
\newpage

\section{The Consulting Scenario}

You work remotely as a consultant specializing in \textbf{spatial intelligence}. Two clients contact you with unrelated problems. In both cases, the client already has information about the situation, but location and spatial relationships may reveal patterns that are difficult to evaluate from ordinary tables alone.

Your role is not to make the final decision for either client. Your role is to transform geographic data into information that helps each client understand the problem and evaluate possible responses.

The two requests illustrate two different contributions of spatial analytics:

\begin{itemize}[leftmargin=*]
    \item \textbf{Spatial constraints and location}: Which places satisfy or violate a geographic rule, and where are feasible alternatives?
    \item \textbf{Spatial dependence}: Does what happens nearby provide information that changes how an observed value should be interpreted?
\end{itemize}

\begin{validationbox}
A map is not yet a spatial analysis. Spatial intelligence begins when location, distance, neighborhood, or geographic configuration becomes part of the analytical question.
\end{validationbox}


\section{Question 1: Which Gas Stations Are Affected, and Where Could They Relocate?}

\begin{executivebox}
A municipal council has modified a rule affecting gas stations. Under the workshop scenario, a gas station cannot operate within 100 meters of a school.

The municipality needs to identify which existing stations are now in a prohibited location. It also needs an initial screening of places where those affected stations could potentially relocate.

The municipality does not need a general map of schools and gas stations. It needs a reproducible way to translate the distance rule into geographic evidence and then narrow the search for alternatives.
\end{executivebox}

\subsection{Conventional planning strategies}

Municipal planners would normally review permits, parcel maps, zoning information, school locations, road access, inspections, and the legal text of the regulation. Local staff may also know which stations are close to schools and which corridors could accommodate relocation.

These sources remain essential. They contain legal, engineering, commercial, and neighborhood information that cannot be recovered from geometry alone.

However, manually checking many locations against a distance rule is difficult to reproduce and becomes increasingly complicated when several spatial restrictions must be considered at the same time.

\subsection{How spatial analysis may help}

Spatial analysis can translate the regulation into geometry. A distance threshold becomes a buffer around every school. Gas stations can then be tested against those protected areas.

The same logic can be reversed when looking for alternatives. Instead of asking which places are prohibited, we can construct an \textbf{opportunity set}: locations that remain after the spatial restrictions included in the model have been removed.

For this workshop, we use the following strategy:

\begin{enumerate}[leftmargin=*]

    \item Load the municipal boundary, schools, and gas stations.

    \item Verify the coordinate reference system and work in a projected CRS so that distance is measured in meters.

    \item Compute the distance from each gas station to the nearest school.

    \item Convert the legal distance rule into 100-meter school buffers.

    \item Identify gas stations whose geometries intersect the protected areas.

    \item Retrieve the drivable road network and retain road types considered plausible for the workshop relocation exercise.

    \item Construct exclusion areas around schools and existing gas stations.

    \item Remove the excluded areas from the suitable road network.

    \item For each affected station, identify several nearby candidate points from the remaining road network.

    \item Present the compliance list and candidate locations to the municipality for further legal, planning, and engineering review.

\end{enumerate}

\begin{notebookbox}
\textbf{Notebook:} \texttt{Spatial\_Q1\_Gas\_Stations\_and\_Schools.ipynb}\\
\textbf{Data:} Boston municipal geometry and OpenStreetMap schools, gas stations, and roads\\
\textbf{Main outputs:} Affected-station table, protected-area map, and screened relocation candidates\\[0.5em]
\colablink
\end{notebookbox}

\outputplaceholder{Gas stations affected by the school-distance rule}

\outputplaceholder{Screened relocation candidates}

\subsection{How to read the outputs}

The first output answers a compliance question: which stations intersect the protected areas created by the distance rule?

The second output answers a different question: which locations remain plausible after applying the spatial constraints included in the notebook?

The reader should ask:

\begin{enumerate}[leftmargin=*]
    \item Which stations are clearly inside the prohibited area?
    \item Are any results sensitive to how point and polygon gas-station geometries are represented?
    \item How does changing the school-distance threshold alter the number of affected stations?
    \item How strongly do the road-type and competition-distance assumptions reduce the candidate set?
    \item Which non-spatial information must still be checked before a location could be considered feasible?
\end{enumerate}

\begin{validationbox}
A candidate location is not an approved site. The notebook does not include zoning, parcel ownership, environmental review, utilities, traffic engineering, construction cost, neighborhood effects, or the full legal interpretation of the regulation.
\end{validationbox}

\begin{decisionbox}
The analysis supports spatial screening. It identifies stations affected by the rule and narrows the geographic search for possible alternatives. The municipality remains responsible for determining whether any candidate location is legally, technically, economically, and politically feasible.
\end{decisionbox}


\section{Question 2: What Kind of Attention May Be Needed Across the Epidemic Map?}

\begin{executivebox}
A public-health organization is monitoring COVID-19 fatality across Peruvian provinces. A choropleth map shows where fatality is currently high and low, but this is not enough to distinguish different spatial situations.

A province with high fatality may be part of a broader high-fatality region or may be an unusual local outlier. A province with relatively low fatality may be located inside a low-fatality environment or may be surrounded by provinces with much higher values.

The client therefore asks a broader question: given the current spatial pattern of fatality, where might different kinds of attention be warranted?
\end{executivebox}

\subsection{Conventional public-health strategies}

Public-health teams normally examine incidence, fatality, hospital capacity, vaccination, mobility, demographic vulnerability, comorbidities, testing, health-system access, and recent trends. Field reports and epidemiological expertise remain essential for interpreting why one place differs from another.

A choropleth map is also useful because it makes geographic variation visible.

However, a choropleth evaluates each area mainly through its own value. It does not formally ask whether the values observed in neighboring areas change how a province should be interpreted.

\subsection{How spatial analysis may help}

Spatial dependence begins with a simple idea: \textbf{nearby places may carry relevant information}.

We first define which provinces count as neighbors and calculate the spatial lag---the average fatality rate among those neighbors. Global Moran's \(I\) then evaluates whether fatality shows spatial dependence across the study area.

Local Moran's \(I\), or LISA, moves from the general pattern to individual provinces. It distinguishes four significant local configurations:

\begin{table}[!htbp]
\centering
\caption{LISA patterns and spatial-intelligence interpretation}
\label{tab:lisa_patterns}
\begin{tabularx}{\textwidth}{lXX}
\toprule
Pattern & Spatial configuration & Possible attention type \\
\midrule

\textbf{High--High (HH)} &
High fatality surrounded by high-fatality neighbors. &
Priority intervention: the area is already part of a broader high-fatality environment. \\

\textbf{High--Low (HL)} &
High fatality surrounded by lower-fatality neighbors. &
Local investigation: the province is an unusual high-value spatial outlier. \\

\textbf{Low--High (LH)} &
Lower fatality surrounded by high-fatality neighbors. &
Early warning / prevention: current local conditions are comparatively better, but the surrounding spatial context is unfavorable. \\

\textbf{Low--Low (LL)} &
Lower fatality surrounded by lower-fatality neighbors. &
Preserve and monitor: both the local and surrounding situation are comparatively favorable. \\

\bottomrule
\end{tabularx}
\end{table}

For this workshop, we use the following strategy:

\begin{enumerate}[leftmargin=*]

    \item Load the provincial COVID-19 spatial data and calculate the 2021 fatality rate.

    \item Begin with a choropleth to understand the geographic distribution of the outcome.

    \item Define neighboring provinces using rook contiguity.

    \item Row-standardize the spatial weights and calculate each province's spatial lag.

    \item Estimate Global Moran's \(I\) to determine whether the study area exhibits spatial dependence.

    \item Estimate Local Moran's \(I\) to classify significant local patterns as HH, HL, LH, or LL.

    \item Translate those patterns into different types of attention for the public-health team.

    \item Compare the classifications with a second neighborhood definition, three-nearest neighbors, to identify conclusions that are sensitive to the spatial assumptions.

    \item Present the spatial patterns to public-health experts for epidemiological interpretation.

\end{enumerate}

\begin{notebookbox}
\textbf{Notebook:} \texttt{Spatial\_Q2\_COVID\_LISA\_Strategy.ipynb}\\
\textbf{Data:} Peruvian provincial COVID-19 data\\
\textbf{Main outputs:} Fatality choropleth, Moran's \(I\), LISA map, spatial-intelligence table, and neighborhood-sensitivity comparison\\[0.5em]
\colablink
\end{notebookbox}

\outputplaceholder{Local spatial patterns of COVID-19 fatality}

\outputplaceholder{Province-level spatial-intelligence table}

\subsection{How to read the outputs}

The categories should not be read simply as ``good'' and ``bad'' areas. They separate different geographic situations.

A High--High province is already located inside a broader high-fatality environment. A High--Low province is different: its high value is unusual relative to its neighbors and may therefore require a more local explanation.

The Low--High category is particularly useful for early-warning thinking. Its current fatality rate is comparatively lower, but the surrounding areas have higher values. This does not mean that fatality will necessarily increase. It means that the province occupies a spatial configuration that deserves attention when combined with epidemiological evidence.

Low--Low provinces provide a different analytical opportunity. The public-health team may want to understand whether favorable conditions are stable and whether any local practices or capacities are contributing to the pattern.

The reader should ask:

\begin{enumerate}[leftmargin=*]
    \item Where are the significant HH, HL, LH, and LL patterns?
    \item Which LH provinces remain classified as early-warning cases under an alternative neighborhood definition?
    \item Which HL provinces require a local explanation because their surrounding context is comparatively favorable?
    \item Are HH areas also places with limited health-system capacity or other known vulnerabilities?
    \item What additional epidemiological evidence is required before an intervention is proposed?
\end{enumerate}

\begin{validationbox}
LISA identifies spatial association, not causation. A Low--High province is not a forecast that fatality will rise. A High--High cluster does not prove transmission between neighboring provinces. Mobility, health-system capacity, demographic composition, data quality, timing, and many other processes may generate the observed pattern.
\end{validationbox}

\begin{decisionbox}
The analysis supports differentiated territorial attention. It helps the public-health team distinguish broad high-fatality clusters, unusual local problems, potential early-warning configurations, and comparatively favorable areas. Public-health experts remain responsible for determining which epidemiological response, if any, is appropriate.
\end{decisionbox}


\section{From Spatial Computation to Client Intelligence}

The two requests use different spatial methods but follow the same consulting logic:

\begin{table}[!htbp]
\centering
\caption{From client request to spatial-intelligence output}
\begin{tabularx}{\textwidth}{>{\raggedright\arraybackslash}p{0.24\textwidth} >{\raggedright\arraybackslash}p{0.31\textwidth} X}
\toprule
Client question & Main spatial operation & Decision-support output \\
\midrule

Which gas stations are affected, and where could they relocate? &
Distance, buffers, spatial joins, overlay, and spatial constraints. &
Compliance screening and candidate-location set. \\

What kind of attention may be needed across the epidemic map? &
Spatial weights, spatial lag, Moran's \(I\), and LISA. &
Territorial classification for differentiated expert attention. \\

\bottomrule
\end{tabularx}
\end{table}

In both cases, the computational method reduces a complex geographic problem to a reproducible set of spatial relationships. The result is not an automated decision. It is evidence that can be combined with domain expertise, institutional rules, and information that is not contained in the spatial data.

\end{document}
